% sec:geo-constants — Constants of Motion
\section{Constants of Motion}
\label{sec:geo-constants}

A generic system of eight first-order ODEs requires eight initial conditions
and admits no simplification.  Symmetries change this: every continuous
symmetry of the metric produces a conserved quantity that eliminates one
degree of freedom from the equations of motion.  The Schwarzschild and
Kerr spacetimes each have enough symmetry to reduce the geodesic problem
from eight equations to a manageable core --- and the Kerr spacetime
possesses an additional hidden symmetry, unrelated to any continuous
isometry, that provides a fourth constant of motion and makes the system
fully separable.

\paragraph{Killing vectors and conserved momenta.}

The connection between symmetry and conservation was established in
\cref{sec:geo-formulations}: if the metric is independent of a coordinate
$x^\alpha$, the corresponding momentum component $p_\alpha$ is conserved
along any geodesic (\cref{eq:geo-pdot-cov}).  The geometric object
underlying this statement is a Killing vector $\xi^\mu$ --- a vector field
satisfying the Killing equation $\cd{\mu}\xi_\nu + \cd{\nu}\xi_\mu = 0$.
The quantity
%
\begin{equation}
  Q_\xi = p_\mu\,\xi^\mu
  \label{eq:killing-conserved}
\end{equation}
%
is then conserved along any geodesic: $\d Q_\xi / \d\lambda = 0$.  To
verify this, apply the Leibniz rule along the geodesic:
$\d Q_\xi/\d\lambda =
  (\dot{x}^\nu\cd{\nu}p_\mu)\,\xi^\mu
  + p_\mu\,\dot{x}^\nu\cd{\nu}\xi^\mu$.
The first term vanishes because $p_\mu$ is parallel-transported along the
geodesic.  For the second, writing $p_\mu = g_{\mu\alpha}\dot{x}^\alpha$
yields the contraction $\dot{x}^\alpha\dot{x}^\nu\,\cd{\nu}\xi_\alpha$,
which is zero because $\dot{x}^\alpha\dot{x}^\nu$ is symmetric while
$\cd{\nu}\xi_\alpha$ is antisymmetric by the Killing equation.
\physnote{When $\xi^\mu = \partial_\alpha$ is a coordinate basis vector,
$Q_\xi = p_\alpha$ --- the Killing conservation reduces to the statement
that $p_\alpha$ is constant when $\pd{\alpha}\,\metric = 0$, exactly
as in \cref{eq:geo-pdot-cov}.}

\paragraph{Energy and angular momentum.}

Both the Schwarzschild and Kerr metrics are stationary (independent of $t$)
and axisymmetric (independent of $\varphi$), yielding two Killing vectors:
$\xi^\mu = (\partial_t)^\mu$ and $\psi^\mu = (\partial_\varphi)^\mu$.  The
associated conserved quantities are the specific energy
%
\begin{equation}
  E = -p_\mu\,\xi^\mu = -p_t
  \label{eq:geo-energy}
\end{equation}
%
and the specific angular momentum (about the symmetry axis)
%
\begin{equation}
  L = p_\mu\,\psi^\mu = p_\varphi \,.
  \label{eq:geo-angmom}
\end{equation}
%
The minus sign in \cref{eq:geo-energy} is a convention that makes $E > 0$
for future-directed geodesics: since $\xi^\mu$ is timelike outside the
ergosphere, $p_\mu\,\xi^\mu < 0$ for a future-directed particle, and the
minus sign compensates.  These are the same quantities derived in
\cref{sec:schw-orbits} for the Schwarzschild geometry
(\cref{eq:schw-energy,eq:schw-angmom}) and used in the Kerr ISCO analysis
(\cref{eq:kerr-circ-E,eq:kerr-circ-L}); the Killing-vector framework makes
their origin and universality explicit.
\xrefnote{Inside the ergosphere, $\xi^\mu$ becomes spacelike and $E$ can be
negative for certain orbits.  This is the mechanism behind the Penrose
process (\cref{sec:kerr-horizons}).}

For the spherically symmetric Schwarzschild spacetime, every geodesic can be
rotated into the equatorial plane by a symmetry transformation, so $E$ and
$L$ together with the norm constraint $\mathcal{H} = \text{const}$ suffice
to reduce the problem to a single ODE for $r(\lambda)$.  The Kerr spacetime
is only axisymmetric, so equatorial restriction is no longer guaranteed ---
geodesics generically move in three spatial dimensions, and two conserved
quantities are not enough.

\paragraph{The Carter constant.}

In 1968, Carter discovered that the Kerr geodesic equations are fully
separable despite the spacetime possessing only two Killing vectors
\cite{Carter1968}.  The additional constant of motion arises from a
rank-two Killing tensor $K_{\mu\nu}$ --- a symmetric tensor satisfying
%
\begin{equation}
  \cd{(\alpha} K_{\mu\nu)} = 0 \,,
  \label{eq:killing-tensor}
\end{equation}
%
where the parentheses denote complete symmetrisation over all three indices.
The Carter constant is
%
\begin{equation}
  \mathcal{Q} = K_{\mu\nu}\,\dot{x}^\mu\,\dot{x}^\nu
    = p_\theta^2 + \cos^2\!\theta
      \left(a^2(-\kappa - E^2) + \frac{L^2}{\sin^2\!\theta}\right) ,
  \label{eq:carter-constant}
\end{equation}
%
where $\kappa = -1$ for timelike and $\kappa = 0$ for null geodesics.  The
first equality is the geometric definition; the second is its explicit
evaluation in Boyer--Lindquist coordinates, where the Killing tensor takes
the form
$K_{\mu\nu} = \Sigma\,l_{(\mu}n_{\nu)} + r^2\,\metric$ with $l^\mu$ and
$n^\mu$ the principal null directions of the Kerr geometry
\cite{Chandrasekhar1983}.
\histnote{Carter found the fourth constant by separating the
Hamilton--Jacobi equation in Boyer--Lindquist coordinates, a decade
before Killing tensors and hidden symmetries were systematically
classified.}

The Carter constant has no analogue in the Schwarzschild geometry, where
$K_{\mu\nu}\,\dot{x}^\mu\dot{x}^\nu$ reduces to
$L_x^2 + L_y^2 + L_z^2$ --- the square of the total angular momentum, which is already determined by the three rotational
Killing vectors of spherical symmetry.  For Kerr, where only axial symmetry
survives, $\mathcal{Q}$ provides genuinely independent information: it
constrains the polar motion in a way that neither $E$ nor $L$ can.

\paragraph{Counting degrees of freedom.}

A geodesic in four-dimensional spacetime has eight phase-space variables
$(x^\mu, p_\mu)$.  Each conserved quantity reduces the effective
dimensionality by one:
%
\begin{enumerate}
\item $E = -p_t$ eliminates $p_t$.
\item $L = p_\varphi$ eliminates $p_\varphi$.
\item $\mathcal{H} = \text{const}$ provides a constraint relating the
  remaining momenta to the coordinates.
\item $\mathcal{Q} = \text{const}$ (Kerr only) constrains the polar
  momentum $p_\theta$.
\end{enumerate}
%
For Schwarzschild, the three constants ($E$, $L$, $\mathcal{H}$) reduce the
system to one effective degree of freedom: the radial equation
$\dot{r}^2 = E^2 - V_{\text{eff}}(r)$ of \cref{sec:schw-orbits}.  For
Kerr, the four constants reduce the system to two first-order equations ---
one for $r(\lambda)$ and one for $\theta(\lambda)$ --- each involving only
its own coordinate and the constants of motion.  This is the separated form
used in analytic studies \cite{Chandrasekhar1983}.
\warnnote{The codebase does not use the separated equations.  It integrates
the full eight-component Hamiltonian system
\cref{eq:geo-xdot,eq:geo-pdot-cov} directly, using adaptive step control
to handle the stiffness near the photon sphere.  The conserved quantities
$E$, $L$, and $\mathcal{H}$ serve instead as accuracy
diagnostics: their drift during integration measures accumulated numerical
error (\cref{sec:geo-conservation}).}

\paragraph{The role of conserved quantities in ray tracing.}

For the full Hamiltonian integration used by the codebase, the conserved
quantities serve two purposes.  First, they provide initial-condition
constraints: the photon's initial momentum must satisfy $\mathcal{H} = 0$,
which fixes one component of $p_\mu$ given the other three and the
initial position.  Second, they provide monitoring quantities: computing $E$,
$L$, and $\mathcal{H}$ at each step and comparing with their initial
values reveals discretisation error without any additional ODE solves.  A
typical well-resolved ray trace maintains $\abs{\Delta E/E} \lesssim \epsilon_a$
and $\abs{\Delta\mathcal{H}} \lesssim \epsilon_a$ throughout.
\physnote{For null geodesics, $E$ can be set to unity by the affine
rescaling freedom (\cref{sec:sr-causal}).  The impact parameters
$b = L/E$ and $q = \sqrt{\mathcal{Q}}/E$ then parametrise the photon
trajectory completely --- these are the coordinates on the observer's
sky that the ray tracer maps to pixel positions.}
